\documentclass[11pt]{article}

\usepackage[utf8]{inputenc}
\usepackage[T1]{fontenc}
\usepackage{amsmath,amssymb}
\usepackage{graphicx}
\usepackage{booktabs}
\usepackage{siunitx}
\usepackage[hidelinks]{hyperref}
\usepackage[round]{natbib}
\usepackage{geometry}
\geometry{margin=1in}

\graphicspath{{figures/}}
% Compile-safe figure include: shows a placeholder box if the image is absent,
% so the paper builds before the figures are copied into paper/figures/.
\newcommand{\safefig}[2][width=0.8\linewidth]{%
  \IfFileExists{figures/#2}{\includegraphics[#1]{#2}}%
  {\fbox{\parbox[c][3cm][c]{0.8\linewidth}{\centering\ttfamily figure placeholder: figures/#2}}}}

\title{Can Explainability Detect Shortcut Learning?\\
\large A Controlled Dose--Response Study of Grad-CAM, Integrated Gradients,\\ and Occlusion Sensitivity on a Medical Image Classifier}
\author{Aanya Mittal}
\date{\today}

\begin{document}
\maketitle

\begin{abstract}
% Grounded in the actual experiment. Revise wording freely; keep the numbers.
Deep networks for medical imaging often rely on \emph{shortcuts}---spurious cues correlated
with the label---rather than disease-relevant features. We ask whether three widely used
post-hoc attribution methods (Grad-CAM, Integrated Gradients, and Occlusion Sensitivity) can
detect such reliance under controlled conditions. We inject a small artificial marker at a
random location into PneumoniaMNIST images, correlating it with the disease label at four
strengths ($\rho \in \{0.0, 0.7, 0.9, 1.0\}$), train twelve classifiers (four strengths,
three seeds), and verify behavioral reliance before trusting any explanation. We measure
attribution mass above a random-box chance level ($\mathrm{AM}^{\star}$), pointing-game hit
rate, and a detection rule that requires both above-chance localization and measured behavioral
reliance. Integrated Gradients and Occlusion Sensitivity localize the marker in all nine
shortcut-reliant models and in none of the three controls, with $\mathrm{AM}^{\star}$ rising
monotonically with $\rho$ (Spearman $+0.95$ and $+0.89$ against $\rho$). Grad-CAM fails to
localize the marker in all twelve models; however, we show this outcome is confounded by the
coarse spatial resolution of its feature map on our architecture and cannot be read as evidence
that Grad-CAM is blind to shortcuts. We also observe that Integrated Gradients places
above-chance mass on the marker even in control models that do not rely on it, underscoring that
attribution alone is not evidence of reliance. These are descriptive findings from three seeds
on a single binary task and are not claimed to be statistically powered or to generalize.
\end{abstract}

% ===========================================================================
\section{Introduction}
\label{sec:intro}
% DRAFT --- to expand in the next iteration.
Deep neural networks frequently achieve high benchmark accuracy by exploiting shortcuts:
decision rules that work on the training distribution but fail under shift
\citep{geirhos2020shortcut}. In medical imaging this is not hypothetical: classifiers have been
shown to key on hospital-specific tokens \citep{zech2018variable}, acquisition confounders
\citep{degrave2021ai}, and surgical skin markings \citep{winkler2019association} rather than
pathology. Because explanation (XAI) methods are routinely proposed as a safeguard, a practical
question follows: \emph{can attribution methods actually reveal when a model relies on a
shortcut?} This work studies that question in a controlled setting where the shortcut is known
by construction, so that detection can be scored against ground truth.

Our contributions are: (i) a controlled, reproducible dose--response benchmark that varies
shortcut strength rather than treating the shortcut as merely present/absent; (ii) a detection
protocol that couples attribution localization with a measured behavioral check, guarding against
false positives; and (iii) a head-to-head comparison of Grad-CAM, Integrated Gradients, and
Occlusion Sensitivity, with an explicit analysis of the resolution confound that affects Grad-CAM.

% ===========================================================================
\section{Related Work}
\label{sec:related}
% DRAFT --- to expand.
\textbf{Shortcut learning and spurious correlations.} \citet{geirhos2020shortcut} frame shortcut
learning as a unifying account of generalization failures. In medical imaging,
\citet{zech2018variable}, \citet{degrave2021ai}, and \citet{winkler2019association} document
concrete shortcuts. \citet{lapuschkin2019unmasking} use explanations to expose ``Clever Hans''
behavior at the dataset level.

\textbf{Attribution methods.} We study three established methods: Grad-CAM
\citep{selvaraju2017gradcam}, Integrated Gradients \citep{sundararajan2017axiomatic}, and
Occlusion Sensitivity \citep{zeiler2014visualizing}.

\textbf{Reliability of attributions.} \citet{adebayo2018sanity} show several saliency methods
fail basic sanity checks; \citet{arun2021assessing} and \citet{saporta2022benchmarking} find
saliency maps are weak localizers in medical imaging. Most directly,
\citet{adebayo2020debugging} and \citet{adebayo2022posthoc} ask whether post-hoc explanations
detect injected spurious signals, reporting that detection succeeds for visible/known artifacts
but fails for unknown or non-visible ones---the nuance our study operates within. Explanation-guided
training \citep{ross2017right,rieger2020interpretations} assumes the shortcut is already known;
we instead ask whether it can be \emph{detected}.

% ===========================================================================
\section{Background}
\label{sec:background}
% DRAFT --- brief definitions of the three attribution methods to be expanded.
Let $f$ be a trained classifier and $x$ an input image. An attribution method produces a map
$A(x)$ assigning an importance value to each pixel for a chosen target class. Grad-CAM weights the
activations of the final convolutional layer by their gradients; Integrated Gradients integrates
input gradients along a straight path from a baseline to $x$; Occlusion Sensitivity measures the
change in the target score as a sliding window masks regions of $x$. Full definitions are deferred
to the expanded version.

% ===========================================================================
\section{Research Questions and Hypotheses}
\label{sec:rq}
\begin{description}
  \item[RQ1 (Localization).] Do attribution maps place above-chance mass on the shortcut marker?
  \item[RQ2 (Detection).] Does a localization+behavior rule separate shortcut-reliant models from controls?
  \item[RQ3 (Consistency).] Do the three methods agree?
  \item[RQ4 (Robustness).] How does localization behave as shortcut strength $\rho$ varies?
\end{description}

% ===========================================================================
\section{Methodology}
\label{sec:method}

\paragraph{Dataset.} We use PneumoniaMNIST from MedMNIST v2 \citep{yang2023medmnist}, a binary
chest-X-ray task (normal vs.\ pneumonia) at $28\times28$ resolution with the official
train/validation/test split. Images are replicated to three channels and normalized with ImageNet
statistics (the backbone is ImageNet-pretrained; for the from-scratch model this is a fixed affine).

\paragraph{Controlled shortcut injection.} We stamp a solid $4\times4$-pixel high-intensity square
marker at a \emph{random} location (recorded per image as a ground-truth box $B$). Injection is
one-sided and controlled by a strength $\rho\in[0,1]$: a disease image (label $1$) receives a marker
with probability $\rho$; a normal image never does. Thus $\rho=0$ is a clean baseline (no markers)
and $\rho=1$ makes the marker a perfectly predictive cue. All injection decisions and positions are
drawn from seeded RNGs and are fully reproducible.

\paragraph{Models and training.} Each model is a classifier trained with AdamW (lr $10^{-3}$,
batch $128$, up to $10$ epochs, cross-entropy). The reported runs use an ImageNet-pretrained
ResNet-18 \citep{he2016resnet} (inputs upsampled to $64\times64$); a from-scratch SmallCNN is also
available. \textbf{Checkpoint selection uses a validation set drawn from the same biased
distribution as training}, so selection reflects a realistic in-distribution pipeline rather than
secretly rewarding the least shortcut-reliant model (which clean-validation selection would do).
We train $4~\text{strengths}\times 3~\text{seeds}=12$ models.

\paragraph{Test-time conditions.} On the held-out test split we evaluate: \emph{Clean} (no marker),
\emph{Concordant} (marker follows the training rule), and \emph{Conflict} (marker stamped on normal
images, breaking the correlation). Attribution is computed on the Conflict set for the model's
\emph{predicted} class.

\paragraph{Metrics.} Let $A(x)$ be an attribution map and $B$ the marker box.
\begin{align}
\mathrm{AM} &= \frac{\sum_{(i,j)\in B}|A(i,j)|}{\sum_{(i,j)}|A(i,j)|},
& \mathrm{AM}^{\star} &= \mathrm{AM} - \frac{|B|}{|x|},
\end{align}
where $|B|/|x|$ is the chance mass of an equal-area box. We also report the \emph{pointing-game}
hit rate (fraction of images whose attribution peak lies in $B$) and, per image, compare
$\mathrm{AM}$ to the $95$th percentile of $\mathrm{AM}$ over $500$ random equal-area boxes
(\emph{null-beat}). Behaviorally, we report clean accuracy, the accuracy gain
$\Delta\mathrm{Acc}=\mathrm{Acc}_{\text{concordant}}-\mathrm{Acc}_{\text{clean}}$, and the
\emph{adversarial flip rate} (AFR): among normal images classified correctly when clean, the
fraction that flip to ``pneumonia'' once the marker is added.

\paragraph{Detection rule.} A method \emph{detects} the shortcut in a model if it beats the random-box
null for a majority of Conflict images \textbf{and} the model is behaviorally reliant
($\mathrm{AFR}>0.3$). Requiring the behavioral condition guards against calling a model
shortcut-reliant purely because an attribution map responds to the visible marker.

% ===========================================================================
\section{Experimental Setup}
\label{sec:setup}
We use PyTorch \citep{paszke2019pytorch} and Captum \citep{kokhlikyan2020captum}. Integrated
Gradients uses a zero baseline and $50$ steps; Occlusion uses a $4\times4$ window with stride $2$
(matched to the marker). Attribution is evaluated on up to $100$ Conflict images per model. Seeds
$\{0,1,2\}$ control weights, data order, and marker placement; cuDNN is set deterministic. All code,
configurations, and per-model outputs are released; figures and tables below are generated directly
from \texttt{per\_model\_metrics.csv}.

% ===========================================================================
\section{Results}
\label{sec:results}

\paragraph{The injected shortcut induces measurable reliance (Table~\ref{tab:behavior}).}
At $\rho=0$ the models ignore the marker behaviorally (AFR $6.8\%$, $\Delta\mathrm{Acc}\approx 0$).
From $\rho=0.7$ onward the marker dominates behavior: AFR exceeds $90\%$, meaning the model flips
the large majority of healthy images to ``pneumonia'' when the marker is stamped on them. At
$\rho=1.0$ the shortcut also \emph{displaces} genuine learning: clean accuracy collapses from
$85\%$ to $48\%$ (with large seed variance), i.e.\ the model can barely classify unmarked images.

\begin{table}[t]
\centering
\caption{Behavioral reliance by shortcut strength $\rho$ (mean\,$\pm$\,SD over 3 seeds).}
\label{tab:behavior}
\begin{tabular}{lccc}
\toprule
$\rho$ & Clean acc.\ (\%) & $\Delta\mathrm{Acc}$ (\%) & AFR (\%) \\
\midrule
0.0 & $85.0\pm1.9$ & $-0.3\pm0.3$ & $6.8\pm6.1$ \\
0.7 & $90.7\pm1.3$ & $+0.6\pm0.2$ & $90.6\pm3.2$ \\
0.9 & $89.6\pm1.7$ & $+1.7\pm0.6$ & $91.2\pm6.9$ \\
1.0 & $48.0\pm15.2$ & $+51.7\pm15.8$ & $98.7\pm0.4$ \\
\bottomrule
\end{tabular}
\end{table}

\paragraph{IG and Occlusion localize the shortcut; Grad-CAM does not (Table~\ref{tab:xai}).}
Attribution mass above chance ($\mathrm{AM}^{\star}$) for Integrated Gradients and Occlusion rises
monotonically with $\rho$ and is well above zero for all shortcut-reliant models, with pointing-game
hit rates reaching $98\%$ and $90\%$ at $\rho=1.0$. Grad-CAM's $\mathrm{AM}^{\star}$ stays near zero
at every strength ($\leq 1.2\%$) and its peak almost never lands in the marker (hit rate $\approx
3\%$).

\begin{table}[t]
\centering
\caption{Attribution localization by method and strength: $\mathrm{AM}^{\star}$ (mean\,$\pm$\,SD, \%),
pointing-game hit rate (\%), and models flagged by the detection rule (of 3 seeds).}
\label{tab:xai}
\begin{tabular}{lccc ccc ccc}
\toprule
 & \multicolumn{3}{c}{Grad-CAM} & \multicolumn{3}{c}{Integrated Gradients} & \multicolumn{3}{c}{Occlusion} \\
\cmidrule(lr){2-4}\cmidrule(lr){5-7}\cmidrule(lr){8-10}
$\rho$ & $\mathrm{AM}^{\star}$ & hit & det & $\mathrm{AM}^{\star}$ & hit & det & $\mathrm{AM}^{\star}$ & hit & det \\
\midrule
0.0 & $-0.2\pm0.1$ & 2  & 0/3 & $3.0\pm1.5$  & 22 & 0/3 & $0.1\pm0.3$  & 3  & 0/3 \\
0.7 & $1.0\pm0.4$  & 3  & 0/3 & $8.5\pm1.1$  & 75 & 3/3 & $7.8\pm2.0$  & 66 & 3/3 \\
0.9 & $0.9\pm0.5$  & 3  & 0/3 & $11.0\pm1.6$ & 84 & 3/3 & $13.7\pm2.6$ & 84 & 3/3 \\
1.0 & $1.2\pm0.3$  & 3  & 0/3 & $19.2\pm2.9$ & 98 & 3/3 & $17.2\pm6.0$ & 90 & 3/3 \\
\midrule
\multicolumn{1}{l}{\textbf{Total det.}} & & & \textbf{0/12} & & & \textbf{9/12} & & & \textbf{9/12} \\
\bottomrule
\end{tabular}
\end{table}

\begin{figure}[t]
\centering
\safefig{amstar_vs_rho.png}
\caption{Attribution mass above chance ($\mathrm{AM}^{\star}$) vs.\ shortcut strength $\rho$
(mean\,$\pm$\,SD over seeds). IG and Occlusion increase with $\rho$; Grad-CAM stays near zero.}
\label{fig:amstar}
\end{figure}

\paragraph{Localization tracks shortcut strength (RQ4).} Across the 12 models, $\mathrm{AM}^{\star}$
correlates with $\rho$ at Spearman $+0.95$ (IG) and $+0.89$ (Occlusion), and with the behavioral
displacement $\Delta\mathrm{Acc}$ at $+0.97$ and $+0.91$. Because AFR saturates at $\rho\geq0.7$,
within the reliant regime $\mathrm{AM}^{\star}$ tracks $\rho$ more than it tracks AFR
(Figure~\ref{fig:faith}). Grad-CAM shows only a weak, small-magnitude trend ($+0.61$ vs.\ $\rho$).

\begin{figure}[t]
\centering
\safefig{afr_vs_amstar.png}
\caption{Behavioral reliance (AFR) vs.\ localization ($\mathrm{AM}^{\star}$) across all 12 models.}
\label{fig:faith}
\end{figure}

\paragraph{A false-positive caution.} At $\rho=0$, Integrated Gradients already places above-chance
mass on the marker (null-beat $48\%$, hit $22\%$) even though the model does not rely on it
(AFR $6.8\%$). Attribution responding to a visible marker is therefore \emph{not} evidence of
reliance; the behavioral gate in our detection rule is what prevents a false alarm.

% ===========================================================================
\section{Discussion}
\label{sec:discussion}
% DRAFT --- to expand.
Under controlled, \emph{visible}, \emph{known} shortcuts, perturbation- and path-integral-based
attributions (Occlusion, Integrated Gradients) are effective localizers and their signal scales
with shortcut strength, consistent with the positive side of \citet{adebayo2020debugging}. The
false-positive behavior of IG at $\rho=0$ echoes \citet{adebayo2022posthoc}: localization alone can
mislead, which is why we tie detection to a behavioral test.

% ===========================================================================
\section{Limitations}
\label{sec:limitations}
% These are load-bearing; keep them.
\begin{itemize}
  \item \textbf{Grad-CAM's failure is confounded by resolution.} On ResNet-18 with $64\times64$
  inputs the final convolutional feature map is only $2\times2$, which cannot resolve a $4$-pixel
  marker. We therefore \emph{do not} conclude that Grad-CAM is blind to shortcuts; a fair comparison
  requires a finer feature map (e.g.\ the SmallCNN variant, $14\times14$) or higher input resolution.
  \item \textbf{Small sample.} Three seeds and one binary task yield descriptive trends only; we make
  no claims of statistical significance or generalization across datasets/architectures.
  \item \textbf{Visible, known shortcut.} We study the easy regime; non-visible or unknown shortcuts
  (cf.\ \citet{adebayo2022posthoc}) are out of scope.
  \item \textbf{Method hyperparameters.} Occlusion window/stride and the IG baseline are fixed a
  priori; results may shift under other choices.
\end{itemize}

% ===========================================================================
\section{Conclusion}
\label{sec:conclusion}
% DRAFT --- to expand.
In a controlled medical-imaging setting, Integrated Gradients and Occlusion Sensitivity reliably
localize an injected shortcut and their localization scales with shortcut strength, while coupling
attribution with a behavioral check is necessary to avoid false positives. Grad-CAM did not localize
the marker, but we attribute this to feature-map resolution rather than a fundamental limitation,
and identify the resolution ablation as the key next step.

\bibliographystyle{plainnat}
\bibliography{references}

\end{document}
